La corrección de exámenes escritos es un proceso intensivo en tiempo que
carece de plataformas adaptadas al contexto académico español. Las
soluciones existentes, mayoritariamente extranjeras, presentan barreras
de idioma e integración incompleta con los sistemas de gestión académica
nacionales.

Este Trabajo Fin de Grado presenta el diseño y la implementación del
\dfn{backend} del SCDEE (Sistema de Corrección Digital de Exámenes
Escritos), un servicio distribuido \dfn{multi-tenant} que automatiza el
ciclo completo de corrección, desde la ingesta de páginas escaneadas
hasta el archivo definitivo del examen corregido. El sistema integra un
\textit{pipeline} de reconocimiento basado en códigos \acs{qr} y
\acs{ocr} tolerante a errores, sobre una arquitectura modular de trece
aplicaciones Django con una capa de servicios extensible mediante
\dfnpl{plugin} intercambiables (autenticación, motores de \acs{ocr},
reconocedores de zonas).

El \dfn{backend} implementa un control de acceso \dfn{multi-tenant} con
permisos contextuales por asignatura, un mecanismo de asignación
inteligente de exámenes a estudiantes y calificación manual y por
rúbrica. La seguridad de los datos sensibles se garantiza mediante
cifrado a nivel de campo, \textit{hashing} robusto de contraseñas y
marcas de agua dinámicas sobre los \acs{pdf} servidos a los alumnos,
respaldado por un sistema de auditoría inmutable de todas las
operaciones. El sistema ofrece además configuración por organización,
que permite adaptar parámetros del reconocimiento, límites y permisos, y
tareas de mantenimiento automatizadas con notificaciones por correo
electrónico.

La validación del sistema combina una cobertura cercana al 90\,\% en
pruebas unitarias y de integración, pruebas de carga con Locust que
verifican el \cref{rnf:concurrency} en operación normal, y una
evaluación empírica de la calidad del \acs{ocr} que justifica la
elección del motor EasyOCR. La arquitectura queda preparada para
evolucionar hacia despliegues en producción, integración con sistemas
\acs{lms} y federación de identidad institucional.

\palabrasclave{corrección digital de exámenes, multi-tenancy, OCR,
Django, sistemas distribuidos, arquitectura backend}
